\SourcePage{312}{315}
\section{Hydrogen-like energy levels}

Hydrogen, the simplest atom, is the only atom for which the Schrödinger
equation admits an exact solution. The energy levels of hydrogen, and also of
the one-electron ions $\mathrm{He}^+$, $\mathrm{Li}^{++}$, \ldots, are given by
the Bohr formula (36.10):
\begin{equation}
E=-\frac{mZ^2e^4}{2\hbar^2(1+m/M)}\frac1{n^2}.
\tag{68.1}
\end{equation}
Here $Ze$ is the nuclear charge, $M$ the nuclear mass, and $m$ the electron
mass. Notice that the dependence on nuclear mass is very weak.

Equation (68.1) disregards every relativistic effect. In this approximation
hydrogen has the special additional (\emph{accidental}) degeneracy already
discussed in \S~36: at fixed principal quantum number $n$, the energy does not
depend on the orbital angular momentum $l$.

Other atoms too have states with properties resembling those of hydrogen.
These are highly excited states in which one electron has a large principal
quantum number and therefore spends most of its time far from the nucleus. To
some approximation, its motion may be treated as motion in the Coulomb field
of the \emph{ionic core}, whose effective charge is unity. The energy-level
values obtained in this way are nevertheless too inaccurate, and require

\SourcePage{313}{316}
a correction that allows for the departure of the short-range field from a
pure Coulomb field. The form of this correction follows readily from the
following argument.

States with large quantum numbers are semiclassical, so their energy levels
may be found from the Bohr--Sommerfeld quantization rules (48.6). At distances
from the nucleus small compared with the ``orbital radius,'' the departure of
the field from the Coulomb form may formally be represented by changing the
boundary condition imposed on the wave function at $r=0$. This changes the
constant $\gamma$ in the quantization condition for radial motion. Since the
rest of that condition remains unaltered, the resulting energy-level formula
must differ from the hydrogen formula by replacing the radial quantum number,
or equivalently the principal quantum number $n$, with $n+\Delta_l$. Here
$\Delta_l$ is a constant known as the \emph{Rydberg correction}:
\begin{equation}
E=-\frac{me^4}{2\hbar^2}\frac1{(n+\Delta_l)^2}.
\tag{68.2}
\end{equation}

By its definition the Rydberg correction is independent of $n$, but it is of
course a function of the excited electron's azimuthal quantum number $l$
(written as the subscript on $\Delta$), and also of the total atomic angular
momenta $L$ and $S$. At fixed $L$ and $S$, $\Delta_l$ decreases rapidly as $l$
increases. The larger $l$ is, the less time the electron spends near the
nucleus, and the more closely the energy levels must approach the hydrogen
levels\footnote{As an illustration, consider the empirical Rydberg corrections
for highly excited states of helium. The total spin of the helium atom may be
$S=0$ or $1$, while in the states under consideration its total orbital
angular momentum is the same as the angular momentum $l$ of the excited
electron (the second electron is in the $1s$ state). The Rydberg corrections
are, for $S=0$,
\[
\Delta_0=-0.140,\qquad \Delta_1=+0.012,\qquad \Delta_2=-0.0022,
\]
and, for $S=1$,
\[
\Delta_0=-0.296,\qquad \Delta_1=-0.068,\qquad \Delta_2=-0.0029.
\] }.

\ProblemHeading
Find the asymptotic expression, at large distances from the ionic core, for
the wave function of an electron in a hydrogen-like $s$ state.

\SolutionHeading At large distances, where $U=-1/r$ in atomic units, the
required function $\psi(r)$ obeys the Schrödinger equation
\[
\psi''+\frac2r\psi'-\varkappa^2\psi+\frac2r\psi=0,
\]

\SourcePage{314}{317}
where $\varkappa=\sqrt{2|E|}$. Seek a solution of the form
$\psi=\operatorname{const}\cdot r^\nu e^{-\varkappa r}$. On neglecting terms in
the equation that decrease faster than $\psi/r$, one obtains
\[
\psi=\operatorname{const}\cdot r^{1/\varkappa-1}e^{-\varkappa r}.
\]
